%%% Local Variables:
%%% mode: latex
%%% TeX-master: "../main"
%%% End:

\chapter[高动态敏捷类步态专家的独立训练]{高动态敏捷类步态专家的独立训练}
\label{cha:agile-experts}

本章对应"分训"阶段的第三类专家——高动态敏捷类步态专家，阐述跳跃专家与后空翻专家两个专家网络的独立训练。跳跃专家面向高台阶与木桩间隙地形，构建单PPO、PPO+SLIP、SLIP+PPO+VMC三种算法范式并对比；后空翻专家面向特殊动作需求，采用确定性翻转方案训练。二者重点解决"既要高动态又要稳姿态"的奖励冲突问题与高动态动作的收敛稳定性问题。

\section{高动态敏捷类专家库概述}

高动态敏捷类专家库包含两个专家：

（1）\textbf{跳跃专家}：面向高台阶与木桩间隙地形，实现指令控制跳跃（跳高/跳远）及行走过程中的跳跃，训练算法涵盖单PPO、PPO+SLIP、SLIP+PPO+VMC三种范式，控制频率200Hz；

（2）\textbf{后空翻专家}：面向平面地形的特殊动作，实现指令控制后空翻，采用基于有限状态机的确定性翻转方案，控制频率200Hz。

高动态运动的核心难点是\textbf{奖励冲突}：跳跃/空翻需要大幅度的爆发动作，而稳定性要求抑制姿态发散，二者奖励信号相互矛盾。若采用单一奖励函数，策略极易陷入"不敢跳"或"跳后失控"的次优解。为此，本章采用阶段式奖励塑形与融合算法两类手段加以应对。

\section{SLIP模型与VMC融合控制基础}

\subsection{SLIP模型}

弹簧负载倒立摆（Spring-Loaded Inverted Pendulum, SLIP）模型将机器人简化为带无质量弹簧腿的点质量，其质心运动由摆动相与着地相交替构成\cite{blickhan1989,geyer2006}。着地相质心轨迹满足
\begin{equation}
  m\,\ddot{\boldsymbol{p}} = m\boldsymbol{g} + k\,(l_0 - \|\boldsymbol{p} - \boldsymbol{p}_{foot}\|)\,\frac{\boldsymbol{p} - \boldsymbol{p}_{foot}}{\|\boldsymbol{p} - \boldsymbol{p}_{foot}\|},
  \label{eq:slip}
\end{equation}
式中，$m$为等效质量，$\boldsymbol{p}$为质心位置，$\boldsymbol{p}_{foot}$为足端位置，$k$为弹簧刚度，$l_0$为自然腿长。SLIP模型为跳跃的起跳、腾空与落地提供了解析的弹道先验，可用于规划参考质心轨迹或生成前馈动作。

\subsection{虚拟模型控制（VMC）}

虚拟模型控制（Virtual Model Control, VMC）\cite{pratt2001}通过在机器人质心处构造虚拟力/力矩，再经雅可比映射为各关节力矩。设虚拟力为$\boldsymbol{F}_{v}$，则关节力矩为
\begin{equation}
  \boldsymbol{\tau} = \boldsymbol{J}^{\mathrm{T}}\,\boldsymbol{F}_{v},
  \label{eq:vmc}
\end{equation}
式中，$\boldsymbol{J}$为任务空间雅可比矩阵。VMC将"质心运动控制"与"关节执行"解耦，便于叠加平衡约束，本文将其用于落地缓冲与姿态保持的力分配。

\section{跳跃专家}

\subsection{单PPO算法}

$jump\_flat\_ppo$ 采用纯强化学习范式，策略直接由观测映射至关节/轮端动作，所有跳跃阶段行为均由奖励驱动学习。该方法实现简单、不依赖模型，但高维动作空间下收敛慢、易受奖励冲突影响。

\subsection{PPO+SLIP融合算法}

$jump\_flat\_ppo\_slip$ 将SLIP模型的弹道先验引入强化学习：由SLIP规划跳跃参考质心轨迹，PPO策略以轨迹跟踪误差作为奖励反馈学习控制动作。该范式将弹道运动的高层规划交给解析模型，强化学习专注低层跟踪控制，显著降低了奖励冲突与学习难度。

\subsection{SLIP+PPO+VMC融合算法}

$jump\_flat\_ppo\_slip\_vmc$ 进一步引入VMC：SLIP负责规划，VMC负责落地与姿态阶段的力分配与缓冲，PPO负责整体策略学习与任务调度。三种先验的融合使各阶段"各司其职"——起跳由SLIP规划主导、空中由策略调节、落地由VMC缓冲，最大限度消解了高动态动作与稳定控制之间的冲突。

\subsection{行走过程中跳跃}

在行走过程中触发跳跃需要策略同时保持步态基础能力并响应跳跃指令。本文采用课程学习降低初始冲突：移动速度由0.2m/s逐步提升至1.0m/s，在低速下先学会"行走+跳跃"的协调，再逐步应对高速下的强冲突。

\section{后空翻专家}

\subsection{翻转阶段划分与机制}

后空翻（Backflip）分为站立、下蹲、起跳、腾空翻转、落地稳定等阶段，其本质是绕横滚/俯仰轴的360$^\circ$翻转弹道运动。与跳跃相比，后空翻对起跳能量、翻转角动量与落地姿态的要求更严苛，且世界系角速度测量存在低估翻转量的物理问题，直接以"翻转完成"作为奖励判据难以收敛。

\subsection{确定性翻转与前馈控制}

针对后空翻训练的收敛困难，本文采用\textbf{确定性翻转}方案：在翻转阶段以有限状态机（FSM）调度前馈动作（FF），屏蔽策略动作，由纯前馈驱动完成翻转；策略网络专注翻转前后与落地后的站立平衡。翻转完成的物理判据为"飞行态翻过身"（上体法向反转）且翻转进度超过阈值，从而避开测量低估问题。该方案使翻转过程可复现、可控，显著提升了训练稳定性。

\subsection{奖励函数设计}

后空翻奖励采用阶段式塑形：站立阶段奖励平衡保持，下蹲阶段奖励下蹲深度，起跳阶段奖励起跳速度与角度，腾空阶段奖励翻转进度与收腿高度，落地阶段惩罚冲击速度并奖励落地姿态。通过各阶段聚焦不同子目标，从根本上消解"翻转动力"与"稳定平衡"的奖励冲突。

\section{多算法对比分析}

\subsection{跳跃稳定性对比}

【待填入：三种算法在不同跳跃高度/距离下的落地成功率、落地后姿态恢复时间等稳定性指标对比。】

\subsection{收敛速度对比}

【待填入：三种算法达到目标跳跃性能所需训练迭代数/时间对比，建议绘制训练曲线图。】

\subsection{追踪误差对比}

【待填入：三种算法在质心轨迹跟踪误差、速度跟踪误差、角度跟踪误差等方面的定量对比。】

\section{实验结果与分析}

【待填入：后空翻确定性翻转的成功率、翻转完整度、落地稳定性等实验结果，并附翻转过程截图序列与奖励曲线。跳跃任务结果在此汇总，并以多算法对比表（表~\ref{tab:jump-compare}）形式呈现。】

\begin{table}[!htbp]
    \centering
    \bicaption{三种跳跃算法性能对比（示意表格，填入真实数据）}{Performance comparison of three jumping algorithms (illustrative, fill in real data)}
    \resizebox{0.92\textwidth}{!}{
        \begin{tabular}{lccc}
            \toprule[1.5pt]
            指标 & 单PPO & PPO+SLIP & SLIP+PPO+VMC \\
            \midrule[0.75pt]
            跳跃成功率(\%) & 【待填】 & 【待填】 & 【待填】 \\
            收敛迭代数 & 【待填】 & 【待填】 & 【待填】 \\
            质心轨迹误差(m) & 【待填】 & 【待填】 & 【待填】 \\
            落地冲击(m/s) & 【待填】 & 【待填】 & 【待填】 \\
            \bottomrule[1.5pt]
        \end{tabular}
    }
    \label{tab:jump-compare}
\end{table}

\section{本章小结}

本章阐述了高动态敏捷类步态专家（跳跃专家与后空翻专家）的独立训练方法。针对跳跃任务，构建了单PPO、PPO+SLIP、SLIP+PPO+VMC三种算法范式并给出系统对比框架；针对后空翻任务，提出了基于FSM的确定性翻转与前馈控制方案，通过阶段式奖励塑形消解奖励冲突。二者作为调度器在高动态地形场景的专家选项，为统一系统集成提供了敏捷运动能力。
